\documentclass[11pt,letterpaper]{article}

\usepackage[margin=1in]{geometry}
\usepackage{amsmath,amssymb}
\usepackage{array}
\usepackage{booktabs}
\usepackage{tabularx}
\usepackage{enumitem}
\usepackage[hidelinks]{hyperref}

\setlength{\parindent}{0pt}
\setlength{\parskip}{0.4em}
\setlist{itemsep=0.15em,topsep=0.35em,parsep=0pt,partopsep=0pt}
\renewcommand{\arraystretch}{1.25}
\newcommand{\evidencebox}[2]{%
    \noindent\textbf{#1}\par
    \begin{center}
    \fbox{%
        \begin{minipage}[c][1.55in][c]{0.91\linewidth}
        \centering
        \textit{Paste screenshot here}\\[0.45em]
        \small #2
        \end{minipage}%
    }
    \end{center}
    \vspace{0.15em}
}

\begin{document}

\title{\textbf{Foundations of Blockchains}\\
\large Smart-Contract Lab Instructions}
\author{}
\date{Due: \underline{\hspace{2.2in}}}
\maketitle

\textbf{Name:} \underline{\hspace{2.6in}}
\hfill
\textbf{Section:} \underline{\hspace{1.4in}}

\section*{Overview: Choose One Question}

Complete \textbf{one} of the three lab questions below. The purpose of this lab
is not merely to make a contract compile: you must deploy it to Remix's local
test blockchain, interact with it, and provide visible evidence that its main
security behavior works. Read the selected lab's \texttt{README.md} and the
comments around every \texttt{TODO} or multiple-choice decision before editing
the Solidity files.

\textbf{Common Remix workflow}

\begin{enumerate}[leftmargin=0.3in]
    \item Open \href{https://remix.ethereum.org}{remix.ethereum.org} in Chrome
    or Firefox. In File Explorer, use \textbf{Clone} to load the course
    repository, or create a workspace and copy in the files named under your
    selected question.
    \item Open the \textbf{Solidity Compiler} tab. Select compiler version
    \textbf{0.8.20 or newer} and compile the main contract together with its
    imports. Resolve every red compiler error before continuing.
    \item Open \textbf{Deploy \& Run Transactions}. Set \textbf{Environment}
    to \textbf{Remix VM}; do not use a public network or real funds. Select the
    named contract, provide any constructor arguments, and click
    \textbf{Deploy}.
    \item Expand the instance under \textbf{Deployed Contracts} and perform
    the functional check specified for your question. If you change the source,
    compile again and deploy a \emph{new} instance: an old deployment still
    contains the old bytecode.
\end{enumerate}

\textbf{Screenshot submission.} Submit clear, labeled screenshots showing
(1) a successful compile with the contract name and compiler version visible,
(2) the contract under \textbf{Deployed Contracts} with its address visible,
and (3) the successful functional check described below. One screenshot may
satisfy more than one item if all required evidence is readable. Keep the
Remix transaction terminal visible when it helps establish that a call
succeeded; a source-code screenshot by itself is not proof of deployment.
Paste the final images into the framed placeholders on your selected question's
submission page. You may leave the other two question pages blank or omit them
from your submitted copy. Confirm that every screenshot is from your final
compiled version and does not present a failed transaction as a successful
result.

\newpage
\subsection*{Question 1: NFT Marketplace Auction}

\textbf{Background.} This question builds an on-chain auction from three pieces:
an ERC-20 payment token, an ERC-721 NFT collection, and a marketplace contract.
The marketplace escrows the seller's NFT, pulls payment tokens from bidders
using allowances, refunds an outbid participant, and settles after the
deadline. The final checkpoint changes the payment rule from a first-price
auction to a second-price auction, where the winner pays the second-highest bid
and receives the excess escrow back.

\textbf{Work to complete.} Follow \path{lab/opt3/README.md}.

Complete \texttt{Q1--Q20} in \path{ERC20.sol}, \path{NFTCollection.sol}, and
\path{Marketplace.sol}, then complete the three code blanks in
\path{SecondPriceMarketplace.sol}. In Remix, open the four files under
\path{lab/opt3/src/} and compile them with version 0.8.20 or newer. Allow Remix to
fetch the OpenZeppelin ERC-721 dependency.

\textbf{Deploy and prepare an auction.} Using the first Remix VM account as
the seller, deploy \texttt{ERC20} with an initial supply, name, and symbol;
deploy \texttt{NFTCollection}; and deploy \texttt{Marketplace} with a market
name. In the NFT contract, call \texttt{mintNFT} and note the new token ID.
Approve the marketplace address to move that NFT. Then call
\texttt{createAuction} with the NFT address, payment-token address, token ID,
a positive starting price, and a Unix end time at least several minutes in the
future.

\textbf{Functional check.} Transfer payment tokens to a second Remix VM
account. Switch to that account, approve the marketplace to spend more than the
intended bid, and call \texttt{bid} with auction index 0 and a price strictly
above the starting price. Call \texttt{getCurrentBidOwner(0)} and
\texttt{getCurrentBid(0)}. The returned owner must be the bidder and the price
must equal the new bid. Also verify with \texttt{ownerOf(tokenId)} that the
marketplace, rather than the seller, holds the NFT while the auction is open.
Your screenshots must show the three deployed contract addresses, the
successful \texttt{createAuction} and \texttt{bid} transactions, and the three
read results proving that the NFT and bid are held in escrow as intended.

\newpage
\section*{Question 1 Screenshot Submission}

Paste one readable image into each frame. Keep addresses and function results
large enough to be checked without zooming into the source code.

\evidencebox{Figure 1A: Compile and deployments}{Show a successful compile and
the deployed addresses of \texttt{ERC20}, \texttt{NFTCollection}, and
\texttt{Marketplace}.}

\evidencebox{Figure 1B: Auction creation and NFT escrow}{Show the successful
\texttt{createAuction} transaction and the \texttt{ownerOf(tokenId)} result
equal to the marketplace address.}

\evidencebox{Figure 1C: Bid and marketplace state}{Show the successful
\texttt{bid} transaction together with \texttt{getCurrentBidOwner(0)} and
\texttt{getCurrentBid(0)} returning the expected bidder and price.}

\newpage
\subsection*{Question 2: Hash-Time-Locked ETH and NFT Exchange}

\textbf{Background.} A hash-time-locked contract (HTLC) locks an asset with a
secret and a deadline. The receiver can claim the asset only by revealing a
secret whose hash matches the stored hash. If the secret is not revealed in
time, the original sender can take the asset back. The ETH contract and NFT
contract are on two different chains. They use the same hash, so revealing
the secret to claim one asset makes the secret available for claiming the
other. The deadlines must still leave enough time for the second claim.

\textbf{Work to complete.} All Question 2 files are in
\path{lab/opt2/remix/}. First complete \texttt{TODO 1--7} in
\path{SimpleHTLC.sol} and compile it with \path{SimpleHTLCChecker.sol}. Next
complete \texttt{Q1--Q8} in \path{SimpleNFTHTLC.sol}, then compile all five
Solidity files in the same folder. Complete the ETH checker before running the
NFT checker because the final scenario imports your completed ETH HTLC as the
ETH contract in the atomic sale.

\textbf{Functional check, ETH contract.} In Remix VM, deploy
\texttt{SimpleHTLCChecker}. Enter \textbf{3} in the \textbf{VALUE} field and
select \textbf{wei}, then call \texttt{runAllTests}. Call
\texttt{isSolved}; it must return \texttt{true}. A revert names the behavior
that remains incorrect. After a fix, recompile and deploy a fresh checker.

\textbf{Functional check, NFT contract and atomic sale.} Deploy
\texttt{SimpleNFTHTLCChecker}. Enter \textbf{1} in the \textbf{VALUE} field
with unit \textbf{wei}, call \texttt{runAllTests}, and then call
\texttt{isSolved}; it must return \texttt{true}. This checker verifies NFT
locking, rejection of an incorrect secret, receiver withdrawal, sender refund
after expiry, and the linked ETH-for-NFT sale. Your screenshots must show both
checkers and both \texttt{isSolved = true} results, with the
successful \texttt{runAllTests} transactions visible.

\newpage
\section*{Question 2 Screenshot Submission}

Paste one readable image into each frame. Include the Remix transaction terminal
and the expanded checker whenever needed to make the result clear.

\evidencebox{Figure 2A: ETH HTLC compile and setup}{Show a successful
compile with compiler version 0.8.20 or newer and the
\texttt{SimpleHTLCChecker} address.}

\evidencebox{Figure 2B: ETH HTLC functional result}{Show the successful
\texttt{runAllTests} transaction with 3 wei attached and
\texttt{isSolved = true}.}

\evidencebox{Figure 2C: NFT HTLC and atomic-sale result}{Show the deployed
\texttt{SimpleNFTHTLCChecker}, the successful \texttt{runAllTests} transaction
with 1 wei attached, and \texttt{isSolved = true}.}

\newpage
\subsection*{Question 3: Smart-Contract Wallet and 2-of-3 Multisig}

\textbf{Background.} An externally owned Ethereum account is controlled by one
private key. A smart-contract wallet replaces that single point of control with
code. In this question you progress from an owner-controlled ETH wallet, through
ERC-20 token handling, to a wallet with three administrators. The final wallet
requires approval from any two of the three administrators before it executes
a payment. This reduces the risk that one lost or compromised key immediately
loses all funds, while introducing the need to encode actions, count distinct
approvals, and prevent replay.

\textbf{Work to complete.} Follow \path{lab/opt1/README.md} and complete the
required choices and code in the three checkpoint folders under
\path{lab/opt1/src/}. For the deployment demonstration, open
\path{lab/opt1/src/2-multisig/MultisigWallet.sol} and its interface. Compile
\texttt{MultisigWallet}, then deploy it with a JSON array containing three
\emph{different} Remix VM account addresses, for example
\texttt{["0x...","0x...","0x..."]}.

\textbf{Functional check.} Send a small amount of test ETH to the deployed
wallet. Encode a call to \texttt{transferEth(recipient, amount)} and use the
same action bytes for every step. From the first administrator account, call
\texttt{approve(action)}. Switch the Remix \textbf{Account} selector to the
second administrator and call \texttt{execute(action)}; that call supplies the
second approval and executes the payment. Verify that the recipient received
the chosen amount and that calling \texttt{execute} again with the same action
reverts. Your functional screenshot must show the two distinct administrator
accounts or their approval events, the successful \texttt{Executed} event or
transaction, and evidence of the resulting balance change. Also capture the
failed replay attempt or the contract's \texttt{executed} value for the action
hash.

\newpage
\section*{Question 3 Screenshot Submission}

Paste one readable image into each frame. A single Remix screenshot may contain
several panels, but each frame must establish the evidence named below.

\evidencebox{Figure 3A: Compile and deployment}{Show a successful compile,
compiler version, the deployed \texttt{MultisigWallet} address, and the three
distinct administrator addresses used for construction.}

\evidencebox{Figure 3B: Two approvals and execution}{Show calls or events from
two different administrator accounts, the successful \texttt{Executed}
transaction, and evidence that the recipient's balance changed.}

\evidencebox{Figure 3C: Replay prevention}{Show that repeating
\texttt{execute(action)} reverts, or show that \texttt{executed} is
\texttt{true} for the action hash after the successful payment.}

\end{document}
